% !TeX root = ../../Book5.tex
% 纲要标题：### 15.1 泛包络代数
% 纲要位置：工作记录/计划纲要.md，第 4549 行。

\section{泛包络代数}
\label{b5:sec:1501}


% 纲要范围：
%
% * 张量代数按 Lie 括号关系取商
% * 泛包络代数的泛性质
% * Lie 代数模与包络代数模
% * 用 Heisenberg 中心商比较多项式代数与 Weyl 代数；增广、余乘法和反自同构解释平凡模、张量模与对偶模
%

一个 Lie 代数表示只要求括号相容，而作用算子及其乘积属于自同态结合代数。泛包络代数以 Lie 代数为生成元，以 Lie 括号为乘法关系，使这些乘积可以在同一个代数中计算。取商构造之后，还须证明原 Lie 代数到泛包络代数的自然映射是单射；这将由 PBW 定理得到。

\subsection{张量代数按 Lie 括号关系取商}
\label{b5:subsec:1501:01}

设 \(\mathfrak g\) 为域 \(k\) 上 Lie 代数。张量代数
\[
	T(\mathfrak g)=\bigoplus_{r\ge0}\mathfrak g^{\otimes r}, \qquad \mathfrak g^{\otimes0}=k,
\]
以张量拼接为乘法，是由向量空间 \(\mathfrak g\) 自由生成的含幺结合代数。给定线性映射 \(\mathfrak g\to A\)，把纯张量送到各因子像的乘积，就唯一得到一个含幺代数同态 \(T(\mathfrak g)\to A\)。
\begin{definition}\label{b5:def:universal-enveloping}
	设 \(\mathfrak g\) 为域 \(k\) 上的 Lie 代数，记 \(T(\mathfrak g)\) 为其底层 \(k\) 向量空间的张量代数。令 \(I\) 为 \(T(\mathfrak g)\) 中由
	\[
		x\otimes y-y\otimes x-[x,y]\qquad(x,y\in\mathfrak g)
	\]
	生成的双边理想，其中 \([x,y]\) 看作次数一的张量。代数
	\[
		U(\mathfrak g)=T(\mathfrak g)/I
	\]
	称为 \(\mathfrak g\) 的\term[fanbaoluo-daishu]{泛包络代数}[universal enveloping algebra]。次数一部分的自然映射记为 \(\iota:\mathfrak g\to U(\mathfrak g)\)。
\end{definition}
\symidx{\(U(\mathfrak g)\)：Lie 代数的泛包络代数}
由关系立即得到 \(\iota\bigl([x,y]\bigr)=\iota(x)\iota(y)-\iota(y)\iota(x)\)， 所以 \(\iota\) 是到 \(U(\mathfrak g)^{-}\) 的 Lie 代数同态，其单射性将在 PBW 定理中证明。将全部正次数张量送到零、常数保持不变的映射零化 \(I\)，故下降为含幺代数同态 \(\varepsilon:U(\mathfrak g)\to k\)，称为增广映射。特别地，\(U(\mathfrak g)\) 的单位非零。

\ProseParagraphBreak

关系含有次数二与次数一的两部分，所以商一般不是由张量次数直接得到的分次代数。不过，“次数不超过 \(r\)”的子空间仍能传到商中，下一节将用这个滤过研究它。

\subsection{泛包络代数的泛性质}
\label{b5:subsec:1501:02}

\begin{theorem}\label{b5:thm:enveloping-universal}
	设 \(\mathfrak g\) 为域 \(k\) 上的 Lie 代数，\(A\) 为含幺结合 \(k\) 代数，\(A^{-}\) 为其交换子 Lie 代数，\(\iota:\mathfrak g\to U(\mathfrak g)\) 为泛包络代数的自然映射。对每个 Lie 代数同态 \(\phi:\mathfrak g\to A^{-}\)，存在唯一含幺 \(k\) 代数同态 \(\widetilde\phi:U(\mathfrak g)\to A\)，使 \(\widetilde\phi\iota=\phi\)。 也就是说，下图在底层 \(k\) 向量空间上交换，虚线额外要求为含幺代数同态，并在这一要求下唯一：
	\[
		\begin{tikzcd}[column sep=large,row sep=large]
			\mathfrak g\arrow[r,"\iota"]\arrow[dr,"\phi"'] &U(\mathfrak g)\arrow[d,dashed,"\widetilde\phi"]\\
			&A
		\end{tikzcd}
	\]
\end{theorem}
\begin{proof}
	先由张量代数的自由性延拓 \(\phi\)。因为 \(\phi\) 保持括号，每个关系元的像为 \(\phi(x)\phi(y)-\phi(y)\phi(x)-\phi\bigl([x,y]\bigr)=0\)。 所以延拓零化双边理想 \(I\)，下降到商。\(U(\mathfrak g)\) 由单位和 \(\iota(\mathfrak g)\) 生成，因此这些生成元上的值唯一决定代数同态。
\end{proof}
这就是泛包络代数的泛性质。它也使构造对同态自然：Lie 代数同态 \(f:\mathfrak g\to\mathfrak l\) 唯一诱导 \(U(f):U(\mathfrak g)\to U(\mathfrak l)\)，复合与恒等都因唯一性而保持。是否把单射映成单射，不能仅从泛性质推出；PBW 将另外给出这一结论。
\begin{proposition}\label{b5:prop:enveloping-quotient}
	设 \(\mathfrak g\) 为域 \(k\) 上的 Lie 代数，\(\mathfrak a\subseteq\mathfrak g\) 为理想，\(\iota:\mathfrak g\to U(\mathfrak g)\) 为泛包络代数的自然映射。令 \(J\) 为 \(U(\mathfrak g)\) 中由 \(\iota(\mathfrak a)\) 生成的双边理想，则
	\[
		U(\mathfrak g/\mathfrak a)\cong U(\mathfrak g)/J.
	\]
\end{proposition}
\begin{proof}
	映射 \(\mathfrak g\to U(\mathfrak g)/J\) 零化 \(\mathfrak a\)，故通过商 Lie 代数分解，再由泛性质延拓为左边到右边的同态。反过来，商映射 \(\mathfrak g\to\mathfrak g/\mathfrak a\) 诱导的包络代数同态零化 \(J\)，得到右边到左边的同态。两复合都在单位及来自 \(\mathfrak g\) 的生成元上为恒等，所以互逆。
\end{proof}
若 \(\mathfrak g\) 交换，定义中的关系只是让两个生成元交换，因而 \(U(\mathfrak g)=\operatorname{Sym}(\mathfrak g)\)。 特别地，一维交换 Lie 代数的包络代数是 \(k[t]\)，其表示可以由任意线性算子给出，不存在类似有限循环群关系 \(T^n=I\) 的限制。

\ProseParagraphBreak

取域 \(k\) 上的三维 Heisenberg Lie 代数 \(\mathfrak h=kx\oplus ky\oplus kz\)，规定 \([x,y]=z\)、\([x,z]=[y,z]=0\)。把基上的括号关系代入包络代数的定义，得到展示
\[
	U(\mathfrak h)\cong k\langle x,y,z\rangle/(xy-yx-z,\ xz-zx,\ yz-zy).
\]
这里各括号表示生成的双边理想，也用 \(x,y,z\) 表示它们在包络代数及其商中的自然像。中心 Lie 理想 \(kz\) 被置零后，\(\mathfrak h/kz\) 是由 \(x+kz,y+kz\) 生成的二维交换 Lie 代数。因此\cref{b5:prop:enveloping-quotient}在这个对象上给出
\[
	U(\mathfrak h)/(z)\cong U(\mathfrak h/kz)\cong k[X,Y].
\]
这个多项式识别可以直接核验：将商中的 \(x,y\) 送到 \(X,Y\)，三条展示关系都为零；反过来，将 \(X,Y\) 送到商中的 \(x,y\)，因两者交换而得到多项式代数同态。两个复合固定单位和各生成元，故互逆。将括号的中心值置零，恰好使原来的两个不交换生成元成为多项式变量。

\ProseParagraphBreak

若改在结合代数中施加 \(z=1\)，则得到第一 Weyl 代数
\[
	U(\mathfrak h)/(z-1)\cong A_1(k)\coloneqq k\langle p,q\rangle/(pq-qp-1).
\]
从左边到右边的同态由 \(x\mapsto p\)、\(y\mapsto q\)、\(z\mapsto1\) 给出；反向同态由 \(p\mapsto x\)、\(q\mapsto y\) 给出，因为左边满足 \(xy-yx=1\)。两复合仍固定生成元，故互逆。在 \(k[t]\) 上令
\[
	D(P)=\frac{\mathrm{d}P}{\mathrm{d}t},\qquad T(P)=tP.
\]
对任意 \(P\in k[t]\)，乘积求导给出 \(D(T(P))-T(D(P))=P\)。所以 \(x\mapsto D\)、\(y\mapsto T\)、\(z\mapsto I\) 给出 \(U(\mathfrak h)/(z-1)\) 在 \(k[t]\) 上的幺性左模；其交换子为非零的恒等算子，说明这个商仍非交换。这里指定的是中心元素的标量作用。由 \(\varepsilon(z-1)=-1\ne0\)，理想 \((z-1)\) 不可能由某个 Lie 理想的自然像生成，因为所有这类自然像及其生成的双边理想都被 \(\varepsilon\) 零化。

\subsection{Lie 代数模与包络代数模}
\label{b5:subsec:1501:03}

\begin{theorem}\label{b5:thm:lie-enveloping-modules}
	设 \(\mathfrak g\) 为域 \(k\) 上的 Lie 代数，\(\iota:\mathfrak g\to U(\mathfrak g)\) 为泛包络代数的自然映射。左 \(\mathfrak g\) 模与幺性左 \(U(\mathfrak g)\) 模给出同一个范畴：在相同 \(k\) 向量空间上，作用通过 \(\iota\) 限制或按泛性质唯一延拓；同态、子模、商模与正合列均保持。
\end{theorem}
\begin{proof}
	一个 Lie 代数表示 \(\rho:\mathfrak g\to\mathfrak{gl}(V)\) 是到结合代数 \(\operatorname{End}_k(V)\) 的交换子 Lie 代数的同态，故由\cref{b5:thm:enveloping-universal}唯一延拓为 \(U(\mathfrak g)\) 的含幺作用。反向限制显然恢复表示条件，唯一性使两种过程互逆。 设 \(V,W\) 的 Lie 代数作用分别为 \(\rho_V,\rho_W\)。一个 \(k\) 线性映射 \(T:V\to W\) 是 Lie 代数模同态，当且仅当
	\[
		T\rho_V(x)=\rho_W(x)T\qquad(x\in\mathfrak g).
	\]
	反复使用此式，得到
	\[
		T\rho_V(x_1)\cdots\rho_V(x_r) =\rho_W(x_1)\cdots\rho_W(x_r)T.
	\]
	这些乘积和单位的线性组合给出整个 \(U(\mathfrak g)\) 的作用，所以 \(T\) 也与延拓后的作用相容。反向限制则立即成立。因此两种同态相同，子空间的不变性也相同，核、商及正合性因而不变。
\end{proof}
包络代数中的乘积由连续施行 Lie 代数算子得到。例如在 \(\mathfrak{sl}_2\) 模上，\(ef\)、\(h^2\) 或它们的和都是自然的算子，但一般不是某个 Lie 代数元素的作用。上一章的 Casimir 算子正需要这种乘积；现在可以把它看成一个与具体表示无关的代数元素。

\ProseParagraphBreak

两个包络代数左模的张量积取在域 \(k\) 上。要在这个空间上作用一个代数元素，需要把它分配到两个因子；要在对偶空间上作用，则需要倒转乘积次序。这两项构造和增广映射都可以由同一个泛性质得到。
\begin{proposition}\label{b5:prop:enveloping-tensor-antipode}
	设 \(\mathfrak g\) 为域 \(k\) 上的 Lie 代数，\(U=U(\mathfrak g)\)，\(\iota:\mathfrak g\to U\) 为自然映射。存在唯一含幺 \(k\) 代数同态
	\[
		\varepsilon:U\longrightarrow k,\qquad \Delta:U\longrightarrow U\otimes_k U,
	\]
	以及唯一 \(k\) 线性反自同构 \(S:U\to U\)，满足
	\[
		\begin{aligned}
			\varepsilon(\iota(x))&=0,\\
			\Delta(\iota(x))&=\iota(x)\otimes1+1\otimes\iota(x),\\
			S(\iota(x))&=-\iota(x) \qquad(x\in\mathfrak g).
		\end{aligned}
	\]
	其中 \(S(1)=1\)、\(S(ab)=S(b)S(a)\)。这些映射满足
	\[
		\begin{gathered}
			(\Delta\otimes\operatorname{id}_U)\Delta =(\operatorname{id}_U\otimes\Delta)\Delta,\qquad S^2=\operatorname{id}_U,\\
			(\varepsilon\otimes\operatorname{id}_U)\Delta =\operatorname{id}_U =(\operatorname{id}_U\otimes\varepsilon)\Delta,
		\end{gathered}
	\]
	最后一行采用 \(k\otimes_k U\cong U\cong U\otimes_k k\) 的自然识别。

	若 \(V,W\) 为幺性左 \(U\) 模，则下列公式分别给出 \(k\)、\(V\otimes_k W\)、\(V^*=\operatorname{Hom}_k(V,k)\) 上的幺性左 \(U\) 模结构：
	\[
		\begin{aligned}
			u\cdot c&=\varepsilon(u)c,\\
			u\cdot(v\otimes w)&=\sum_r a_rv\otimes b_rw \quad\left(\Delta(u)=\sum_r a_r\otimes b_r\right),\\
			(u\cdot\varphi)(v)&=\varphi\bigl(S(u)v\bigr).
		\end{aligned}
	\]
	这里 \(u\in U\)、\(c\in k\)、\(v\in V\)、\(w\in W\)、\(\varphi\in V^*\)；张量的求和为有限和。
\end{proposition}
\symidx{\(\varepsilon,\Delta,S\)：包络代数的增广、余乘法与反自同构}
\begin{proof}
	零 Lie 代数同态 \(\mathfrak g\to k^{-}\) 由\cref{b5:thm:enveloping-universal}唯一延拓为 \(\varepsilon\)，它与前面构造的增广映射相同。对 \(\Delta\)，在 \(U\otimes_k U\) 中，不同因子的元素交换，所以
	\[
		\begin{aligned}
			&[\iota(x)\otimes1+1\otimes\iota(x),\, \iota(y)\otimes1+1\otimes\iota(y)]\\
			&\qquad=\iota([x,y])\otimes1+1\otimes\iota([x,y]).
		\end{aligned}
	\]
	因此指定的线性映射保持 Lie 括号，同一泛性质给出良定义且唯一的 \(\Delta\)。两个三重余乘法都是到 \(U^{\otimes3}\) 的含幺代数同态，并将 \(\iota(x)\) 送到
	\[
		\iota(x)\otimes1\otimes1 +1\otimes\iota(x)\otimes1 +1\otimes1\otimes\iota(x).
	\]
	它们在生成元和单位上一致，故相等。用 \(c\otimes u\mapsto cu\)、\(u\otimes c\mapsto cu\) 作自然识别后，两种含 \(\varepsilon\) 的复合都固定 \(\iota(x)\) 和单位，同样得到余单位等式。

	对 \(S\)，把目标换为反代数 \(U^{\mathrm{op}}\)，其乘法次序与 \(U\) 相反。于是
	\[
		[-\iota(x),-\iota(y)]_{\mathrm{op}} =\iota(y)\iota(x)-\iota(x)\iota(y) =-\iota([x,y]).
	\]
	所以 \(x\mapsto-\iota(x)\) 是到 \((U^{\mathrm{op}})^{-}\) 的 Lie 代数同态，由泛性质唯一延拓为 \(U\to U^{\mathrm{op}}\) 的含幺代数同态，也就是反同态 \(S\)。它在生成元乘积上满足
	\[
		S\bigl(\iota(x_1)\cdots\iota(x_r)\bigr) =(-1)^r\iota(x_r)\cdots\iota(x_1).
	\]
	\(S^2\) 为代数同态，并固定所有生成元和单位，故为恒等。这样 \(S\) 可逆，确为反自同构。

	\(\varepsilon\) 是含幺代数同态，所以第一条作用给出左模。映射
	\[
		U\otimes_k U\longrightarrow\operatorname{End}_k(V\otimes_k W),\qquad a\otimes b\longmapsto\bigl(v\otimes w\mapsto av\otimes bw\bigr)
	\]
	由双线性得到，并保持乘法和单位。它与 \(\Delta\) 复合，给出第二条作用；这也说明该公式不依赖 \(\Delta(u)\) 的具体张量分解。对第三条作用，若 \(u,a\in U\)，则
	\[
		\begin{aligned}
			\bigl(u\cdot(a\cdot\varphi)\bigr)(v) &=\varphi\bigl(S(a)S(u)v\bigr)\\
			&=\varphi\bigl(S(ua)v\bigr) =\bigl((ua)\cdot\varphi\bigr)(v).
		\end{aligned}
	\]
	单位同样作用为恒等，故对偶也得到幺性左模。
\end{proof}
限制到 \(\mathfrak g\) 后，第一条作用使每个 \(x\) 为零，给出平凡模；后两条则恢复\cref{b5:prop:lie-dual-tensor-hom}中的张量和对偶作用。以下把 \(\iota(x)\) 简写为 \(x\)。因为 \(\Delta\) 保持乘法，
\[
	\Delta(x^2)=x^2\otimes1+2x\otimes x+1\otimes x^2,
\]
所以
\[
	x^2(v\otimes w)=x^2v\otimes w+2xv\otimes xw+v\otimes x^2w.
\]
特征二时中间项为零；任意特征下，高次乘积的张量作用都由生成元的作用和 \(\Delta\) 的乘法相容性确定。

\ProseParagraphBreak

对偶上的负号来自 \(S(x)=-x\)：\((x\varphi)(v)=-\varphi(xv)\)，而反同态性质保证继续作用时乘积次序正确。例如 \((xy\varphi)(v)=\varphi(yxv)\)。这个负号也与一次项的取逆相容：在 \(U\otimes_k k[\delta]/(\delta^2)\) 中，\((1+\delta x)^{-1}=1-\delta x\)。增广 \(\varepsilon\)、余乘法 \(\Delta\) 和反自同构 \(S\) 因而分别控制平凡模、张量模与对偶模；\chapref{b5:ch:21}将把它们作为同一套 Hopf 代数结构数据统一讨论。
